\section{Gradient Magnitude Regularization for Model Merging}
\label{sec:gradient_magnitude_regularization}

\subsection{Motivation}

Our analysis of metric importance (Section~\ref{sec:feature_importance_comparison}) reveals that gradient-based metrics---specifically \texttt{encoder\_gradient\_l2\_distance} and \texttt{input\_gradient\_l2\_distance}---are the most consistently selected predictors of mergeability across merging methods. These metrics measure how similarly two fine-tuned models respond to input perturbations, capturing functional similarity beyond weight-space comparisons. This finding motivates a regularization strategy that directly encourages gradient proximity during fine-tuning.

\subsection{Method}

Gradient magnitude regularization introduces a penalty on the norm of the gradient at each optimization step:
\begin{equation}
\mathcal{L}_{\text{total}} = \mathcal{L}_{\text{CE}} + \lambda \sum_l \|\nabla_{\theta^{(l)}} \mathcal{L}_{\text{CE}}\|_2^2
\end{equation}
where $\mathcal{L}_{\text{CE}}$ is the cross-entropy loss, $\theta^{(l)}$ are the parameters of layer $l$, and $\lambda$ controls the regularization strength. By penalizing large gradients, this regularization encourages smoother optimization trajectories with smaller parameter updates. Implicitly, this brings different fine-tuned models closer together in gradient space: models trained with constrained gradients will exhibit more similar gradient responses to shared inputs, reducing the gradient-based distance metrics that our analysis identifies as key mergeability predictors.

\subsection{Experimental Setup}

We evaluate gradient magnitude regularization ($\lambda = 1$) on pairwise model merging using the ViT-B-16 architecture. Pairwise merging combines exactly two fine-tuned models at a time, providing a controlled environment to study merging dynamics. We evaluate all 190 unique pairs from the N20 benchmark (20 datasets), computing the average normalized accuracy across both tasks for each pair.

\subsection{Results}

Table~\ref{tab:pairwise_grad_mag} presents the average normalized accuracy across all 190 pairwise merges for five merging methods.

\begin{table}[h]
\centering
\caption{Effect of gradient magnitude regularization on pairwise merging (N20 benchmark, 190 pairs). $\Delta$ indicates absolute percentage point change from baseline.}
\label{tab:pairwise_grad_mag}
\begin{tabular}{lccc}
\toprule
\textbf{Merger} & \textbf{Baseline} & \textbf{Grad Mag ($\lambda$=1)} & \textbf{$\Delta$} \\
\midrule
Arithmetic & 0.9102 & 0.9174 & +0.72\% \\
DARE       & 0.9563 & 0.9628 & +0.65\% \\
TSV        & 0.9818 & 0.9857 & +0.39\% \\
TIES       & 0.8415 & 0.8404 & $-$0.11\% \\
Weight Avg & 0.9586 & 0.9625 & +0.39\% \\
\midrule
\textbf{Average} & \textbf{0.9297} & \textbf{0.9338} & \textbf{+0.41\%} \\
\bottomrule
\end{tabular}
\end{table}

\subsection{Analysis}

Gradient magnitude regularization improves pairwise merging performance for four out of five merging methods, with an average improvement of +0.41\% in normalized accuracy.

\paragraph{Consistent Improvements for Most Mergers.}
Arithmetic (+0.72\%), DARE (+0.65\%), TSV (+0.39\%), and Weight Averaging (+0.39\%) all benefit from gradient magnitude regularization. This is consistent with our metric importance analysis, which identified gradient-based metrics as the top predictors of mergeability for these methods. By regularizing gradient magnitudes during fine-tuning, we reduce the gradient distance between models, directly optimizing the metrics that most strongly predict successful merging.

\paragraph{Why TIES Does Not Improve.}
TIES is the only merging method that exhibits slight performance degradation ($-$0.11\%) with gradient magnitude regularization. This exception is explained by our metric importance analysis: unlike the other four mergers, gradient metrics are \emph{not} among the most dominant predictors for TIES. Table~\ref{tab:metric_ranks} shows the top-5 metrics by coefficient magnitude for each merger. For Weight Averaging, Arithmetic, TSV, and DARE, gradient metrics (\texttt{encoder\_gradient\_l2\_distance} and \texttt{input\_gradient\_l2\_distance}) occupy positions \#1 and \#2. In contrast, TIES has \texttt{interaction\_matrix\_overlap\_bottom\_k} and \texttt{task\_vector\_l2\_distance} as its top predictors, with gradient metrics appearing only at positions \#3 and \#5. Consequently, a regularization strategy designed to improve gradient proximity does not address the factors most relevant to TIES performance.

\begin{table}[h]
\centering
\caption{Top-5 metrics by L1 LOTO coefficient magnitude for each merger. Gradient metrics (bold) dominate for all mergers except TIES.}
\label{tab:metric_ranks}
\small
\begin{tabular}{lccccc}
\toprule
\textbf{Merger} & \textbf{\#1} & \textbf{\#2} & \textbf{\#3} & \textbf{\#4} & \textbf{\#5} \\
\midrule
Weight Avg & \textbf{enc\_grad} & \textbf{inp\_grad} & act\_cos & rs\_top & rs\_bot \\
Arithmetic & \textbf{inp\_grad} & \textbf{enc\_grad} & tv\_l2 & im\_bot & sv\_overlap \\
TSV & \textbf{enc\_grad} & \textbf{inp\_grad} & rs\_bot & rs\_top & act\_cos \\
TIES & im\_bot & tv\_l2 & \textbf{inp\_grad} & sv\_overlap & \textbf{enc\_grad} \\
DARE & \textbf{enc\_grad} & \textbf{inp\_grad} & act\_cos & rs\_top & rs\_bot \\
\bottomrule
\end{tabular}
\end{table}

\paragraph{Quantifying the Gradient Contribution.}
To make fair comparisons across mergers, we normalize the absolute coefficient values to sum to 1, yielding the relative contribution of each metric. Table~\ref{tab:gradient_contribution} shows the combined contribution of gradient metrics for each merger. For Weight Averaging, Arithmetic, TSV, and DARE, gradient metrics collectively account for 38--49\% of the total predictive signal. In contrast, TIES allocates only 28\% to gradient metrics, instead relying more heavily on \texttt{interaction\_matrix\_overlap\_bottom\_k} (16.8\%) and \texttt{task\_vector\_l2\_distance} (16.4\%). This quantitative difference explains why gradient magnitude regularization improves the former group but not TIES: the regularization targets nearly half of the predictive signal for most mergers, but less than one-third for TIES.

\begin{table}[h]
\centering
\caption{Normalized contribution of gradient metrics to mergeability prediction (ViT-B-16). Values represent the fraction of total predictive weight allocated to each metric after normalizing absolute coefficients to sum to 1.}
\label{tab:gradient_contribution}
\begin{tabular}{lccc}
\toprule
\textbf{Merger} & \textbf{enc\_grad\_l2} & \textbf{inp\_grad\_l2} & \textbf{Total} \\
\midrule
Weight Avg & 24.8\% & 19.3\% & \textbf{44.1\%} \\
Arithmetic & 18.4\% & 19.8\% & \textbf{38.2\%} \\
TSV & 23.0\% & 17.4\% & \textbf{40.3\%} \\
TIES & 13.1\% & 15.3\% & 28.4\% \\
DARE & 27.9\% & 20.6\% & \textbf{48.5\%} \\
\bottomrule
\end{tabular}
\end{table}

This finding validates our metric-driven approach to regularization design: regularization strategies should target the metrics that are most predictive for the specific merging method being used. A universal regularization strategy may not benefit all mergers equally, and the effectiveness of a particular regularization can be predicted by examining which metrics dominate the mergeability prediction for that method.

\subsection{Summary}

Gradient magnitude regularization provides consistent improvements for merging methods where gradient-based metrics are the dominant mergeability predictors (Arithmetic, DARE, TSV, Weight Averaging). The exception of TIES---where gradient metrics are not the top predictors---confirms that effective regularization strategies should be informed by metric importance analysis. This demonstrates a practical application of our mergeability prediction framework: the learned metric coefficients not only predict merging success but also guide the design of regularization strategies for improved fine-tuning.
